\section{Reproducibility}\label{app:repro}

All results are produced by the code in the accompanying repository
(\texttt{github.com/scjoshi/MPME}). From the repository root,
\begin{center}
\texttt{make technote}
\end{center}
regenerates every number, table and figure with the fast (implicit-Jacobian) settings and
compiles this note; \texttt{make technote-reference} reruns the Monte Carlo experiments with
the forward-difference reference settings used for \cref{tab:mc,tab:mcb1,tab:mcdiag,tab:mag}.

\begin{table}[!htbp]
\centering
\caption{Code map.}
\small
\begin{tabular}{ll}
\toprule
Component & File\\
\midrule
Protocol, echo timing & \texttt{src/mpme/sequence.py}\\
Steady state (isochromat, EPG) & \texttt{src/mpme/epg.py}\\
Voxel signal \eqref{eq:signal} & \texttt{src/mpme/signal.py}\\
Published inverse & \texttt{src/mpme/paper.py}\\
Model-fit baseline & \texttt{src/mpme/analytic.py}\\
Complex ML, magnitude LS, branch test & \texttt{src/mpme/mle.py}\\
Exact Jacobians & \texttt{src/mpme/fastjac.py}\\
Fisher information, CRLB, Rician, uncertainty maps & \texttt{src/mpme/crlb.py}\\
Digital phantom, low-resolution scan 2, polynomial fit & \texttt{src/mpme/phantom.py}\\
\midrule
CRLB tables and figures (\cref{sec:information}) & \texttt{scripts/crlb\_analysis.py}\\
Estimator efficiency (\cref{tab:mc,tab:mcb1,tab:mcdiag}) & \texttt{scripts/estimator\_efficiency.py}\\
Magnitude comparison (\cref{tab:mag}) & \texttt{scripts/magnitude\_comparison.py}\\
One start vs four (\cref{sec:ml}) & \texttt{scripts/start\_count.py}\\
Jacobian benchmark (\cref{tab:bench}) & \texttt{scripts/benchmark\_fit.py}\\
Phantom experiment (\cref{sec:phantom}) & \texttt{scripts/phantom\_experiment.py}\\
Figures and tables of this note & \texttt{scripts/technote\_figures.py}\\
Unit tests & \texttt{tests/} (\texttt{pytest})\\
\bottomrule
\end{tabular}
\end{table}

\paragraph{Settings.} Noise seeds are fixed in every script (Monte Carlo: the SNR value; phantom:
seed 1 for noise, seed 0 for the phantom). Software: Python 3.12.2, PyTorch 2.6.0, NumPy 2.3.5,
SciPy 1.16.3, Matplotlib 3.10.7; laptop CPU with 8 threads, no GPU. Run times quoted in the note
refer to this machine.

\paragraph{Assumptions not stated in~\cite{cheng2019}.} The time from the RF pulse to the first
echo (4.5\,ms), instantaneous RF pulses, and the polynomial degree used to smooth the flip-angle
map (4 here).
